\documentclass{article}
\usepackage{amsmath}
\begin{document}
\title{A common release check for lending a live acquisition reserve}
\maketitle

\paragraph{Question and kernel.}
Can a locally available optional effect be allowed now without losing
a mandatory effect at the next qualified opportunity? The decision is
made before the time of that opportunity is observed. This note covers
one optional effect and one mandatory future effect on a declared
resource, with already entitled operations held fixed in both arms.
It does not reduce arbitrary live recovery to one release number.
The two executable interpretations are a nominal camera capture slot
and a receiver's rolling byte budget. No claim equates a frame with
bytes, or predicts arrival, motion, permission, or link reliability
from either quota.

\paragraph{Frozen decision rule.}
Assume a scheduled admissible opportunity at an unknown time in a
nonempty set \(T\), with attained first and last times \(a,b\).
For each of the two current choices \(x\), a source adapter must
justify the exact predicate \(\operatorname{eligible}(x,t)\)
on all \(T\). In this family it has form \(e_x\leq t\),
with \(e_x=\infty\) for permanent refusal. \(D\) is a separate
qualified completion cutoff on the opportunity time, after accounting
for downstream processing. Neither \(T\) nor \(D\) is changed by
the optional effect. The optional effect weakly increases \(e\).
The executable guard authorizes the optional effect if
\(e_{\rm use}\leq a\) and \(b\leq D\); otherwise it withholds it.
It reports the family as outside the guaranteed envelope if
\(e_{\rm keep}> a\) or \(b>D\).
Necessity checks the two attained extreme times, sufficiency
uses order comparison. With hard future service priority and no
value assigned to an unsafe optional effect, this maximizes early
optional service among the two admitted choices. A direct policy
checking both choices at all support times agrees; the guard
has no information about which support time is realized.

\paragraph{Capture adapter.}
The pinned Gazebo3 caller starts an ordinary manager pass on an
unrelated active member even while the target's wake mask is live.
The previous target fire advances its due date to \(d\).
If an ordinary target update at the early wake is nominally due,
running it changes the next date to \(d+P\), \(P=50\) ms.
Otherwise the early attempt is a no-op. The availability pair
is these dates. Withholding a due update requires an explicit
per-effect guard, modeled by filtering the target from that
one manager loop, not by forcing a frame or pausing the whole
render pass. All other stock calls in the selected fixture are
unchanged, including both old sensors' callbacks. This change
is not presented as a switch in the unmodified runtime.
The target remains wake-masked, so its changed due date does not
change either old sensor's opportunities within the tested interval.
Pre-opening target captures are optional in this envelope; if they
also supply an old user, the filter is not justified by this proof.
The far-site physical cutoff uses 105 ms modeled decoding, not
an achieved render latency.

\paragraph{Receiver adapter.}
Before choosing an optional packet, the exact live receiver debit
history and packet lengths are known in the selected test. For a
future packet of length \(n\leq C\), remove chronological debits
until free capacity reaches \(n\); the first eligible time is
the corresponding expiry. For 1024 ticks per second, a debit
at \(u\) is still charged at \(u+1024\) and first absent at
\(u+1025\). Recompute with and without the optional debit.
The actual packet admission/reply in the test is pinned SubT
receive, transmit, dispatch and callback code in a deterministic
link fixture. The old and optional sends fit initially, and
there are no additional transmission charges from failed attempts
or other traffic before the required send. The sender's live
debit is contained in the receiver's accepted debit sum here;
a successful receive-cap check therefore also meets transmit
capacity. The queue-side guard is not evidence of control over
uncooperative senders. With unknown later debits or an authorized
sender change, the scalar release must be re-established; the
two values are not passive long-lived state summaries.

\paragraph{Replay and comparison.}
Each source adapter is applied before a sweep over its conditional
opportunity support. Both use the same guarded decision.
The camera tests use the nonnegative nominal times and
a translated modeled capture opening. The radio tests charge the
same existing accepted report in every arm. Stock work-conserving
execution loses a required effect on some traces where the scope
guard preserves it; some optional effects can still proceed when
they cannot spend the reserved future slot. The direct two-choice
enumeration with the same dispatch power, source information and
future guarantees matches every guard decision. This distinguishes
a scoped controlled-effect change from unmodified execution; it
does not demonstrate an advance over factored dispatch or a
published architecture capable of the same filtering. The choice of
future envelopes and adapter was not independently held out. No
physical hardware, concurrent end-to-end run, or independent mission
is claimed.
\end{document}
